\appendix
\label{cha:appendix}

The appendix is an optional part of the manuscript. It can be used to provide additional information that is not essential to the main text. This can include additional tables, figures, or other supplementary material. The appendix is not meant to be a second main part of the manuscript, but rather a complementary part. Typical use cases for the appendix include long mathematical derivations or bulky figures (e.g., from empirical experiments) that would interrupt the reading flow of the main text.

\section{Artificial intelligence use declaration}
To ensure transparency and uphold academic integrity, the following table outlines the artificial intelligence (AI) tools used during the preparation of a thesis. The table is just an example and should be adapted to the individual use cases. If you did not use any AI tools, delete this appendix section.

\begin{table}[h!]
	\centering
	\caption{AI use declaration}
	\label{tab:ai_use_declaration}
	\begin{tabular}{@{}p{2.5cm}@{\hspace{5mm}}p{7cm}@{\hspace{5mm}}p{3.5cm}@{}}
		\toprule
		\textbf{Tool} & \textbf{Use case}               & \textbf{Scope}    \\ \midrule
		ChatGPT 5.0   & Text generation (initial draft) & Chapter 1         \\
		ChatGPT 5.0   & Grammar and spell checking      & Entire manuscript \\
		DALL-E 3      & Image generation                & Figure 2.1        \\
		Claude 4.5    & Python programming              & Chapter 3.2       \\
		\bottomrule
	\end{tabular}
\end{table}

\section{Further appendix section}
Place supplementary material here when it supports the thesis but would interrupt the main text. Examples include extended derivations, additional result tables, questionnaires, or technical documentation. Remove this example section when it is not needed.
